\chapter{Conclusiones}
\label{sec:conclusiones}

Se diseñó y simuló un circuito acondicionador de señal que adapta la salida del sensor LM335, de \SI{2.73}{\volt} a
\SI{3.23}{\volt}, al rango de \SI{0}{\volt} a \SI{5}{\volt} del conversor A/D para temperaturas entre
\SI{0}{\celsius} y \SI{50}{\celsius}. La configuración restadora con referencia por divisor y seguidor permitió
trabajar con una única alimentación de \SI{5}{\volt}, evitando las tensiones intermedias superiores a \SI{15}{\volt}
de la topología que amplifica antes de restar.

La simulación verificó la ganancia de diseño, con una sensibilidad de $\SI{99.94}{\milli\volt/\celsius}$ frente a los
$\SI{100}{\milli\volt/\celsius}$ requeridos. El error se concentra en un desplazamiento constante de cerca de
\SI{0.5}{\celsius}, debido a que la referencia obtenida con resistencias normalizadas ($\SI{2.727}{\volt}$) es menor
que la tensión real del sensor a \SI{0}{\celsius} ($\SI{2.7315}{\volt}$), y en la saturación de la salida en el último
medio grado del rango.

Ambos efectos pueden corregirse ajustando la referencia: reemplazando parte del divisor por un potenciómetro de ajuste
para llevar $V_{REF}$ a $\SI{2.7315}{\volt}$, o calibrando el sensor mediante su terminal ADJ. También puede
compensarse en el microcontrolador, restando el desplazamiento medido. Con estas consideraciones se cumplen los
objetivos del trabajo práctico.
